\documentclass[11pt]{article}

% ---- Packages (all standard; compiles with pdflatex on TeX Live / MiKTeX / Overleaf) ----
\usepackage[utf8]{inputenc}
\usepackage[T1]{fontenc}
\usepackage[margin=1in]{geometry}
\usepackage{amsmath,amssymb,amsthm}
\usepackage{booktabs}
\usepackage{array}
\usepackage{enumitem}
\usepackage{xcolor}
\usepackage{mathtools}
\usepackage[hidelinks]{hyperref}

\definecolor{navy}{HTML}{0C1E3E}
\definecolor{gold}{HTML}{8B6914}
\definecolor{rust}{HTML}{8B3A0E}
\definecolor{sage}{HTML}{3F5A3B}

\newcommand{\vecx}{\operatorname{vec}}
\newcommand{\Cov}{\operatorname{Cov}}
\newcommand{\CVaR}{\mathrm{CVaR}}
\newcommand{\R}{\mathbb{R}}
\newcommand{\E}{\mathbb{E}}
\newcommand{\gco}{\mathrm{gCO_2eq/kWh}}
\newcommand{\degC}{^{\circ}\mathrm{C}}

\title{\vspace{-2.2em}\textbf{\color{navy} Spatially-Correlated Distributionally Robust\\[2pt]
Scheduling for Carbon-Aware Data Centers}\\[6pt]
\large\color{gold} Status Report --- Revised Feasible Set and a Second-Grid Generalization}
\author{Marco Ortiz Togashi \\ \small Supervisor: Prof.\ Bissan Ghaddar \\
\small IE School of Science \& Technology --- Capstone in Operations Research}
\date{6 June 2026}

\begin{document}
\maketitle
\thispagestyle{empty}

\begin{abstract}
\noindent We report this period's progress on the distributionally robust optimization (DRO) approach
to carbon-aware data-center scheduling. The central question is whether modeling the \emph{spatial
correlation} of carbon intensity across electrically coupled regions improves a robust schedule. We
(i) replaced a previously dominant aggregate power cap with a more defensible feasible set --- adding a
job deadline constraint and a temperature-coupled cooling (PUE) constraint driven by a new external
temperature dataset --- and (ii) transported the \emph{identical} methodology to a second, independent
grid (Iberia and France) to test generalization. Using a falsification test (the schedule fitted to
the empirical joint covariance versus one with all cross-region structure destroyed), evaluated by
out-of-sample Conditional Value-at-Risk on a held-out year under strict pre-registration, we find that
spatial correlation yields \textbf{no robust scheduling value} in either grid: the worst-case tail
emissions change by at most $0.06\%$ across all constraint regimes and flexibility levels. We
demonstrate the null is \emph{substantive} (the optimizer is not frozen) rather than a mechanical
artifact, that it survives a shrinkage and residual-space robustness battery, and --- notably --- that
it holds even though Spain and Portugal have an empirical carbon-intensity correlation of $\approx
0.85$, higher than any California pair. The methodological contribution is a clean, replicated negative
result with a mechanistic explanation.
\end{abstract}

\section{Introduction and scope of this period}
The project extends a Wasserstein-DRO framework for carbon-aware scheduling by treating carbon
intensity as a stochastic vector across multiple co-regional clusters, so that \emph{cross-region}
(spatial) correlation can in principle inform the schedule. This report covers two work items:
\begin{enumerate}[leftmargin=1.4em,itemsep=2pt]
  \item \textbf{A revised feasible set.} The earlier model relied on a single aggregate power cap that
  dominated the geometry of the feasible region. We removed it and rebuilt the operational constraints
  from more physically grounded components --- a deferral-deadline constraint and a temperature-coupled
  Power-Usage-Effectiveness (PUE) constraint --- and re-validated the central experiment across three
  constraint regimes of increasing tightness.
  \item \textbf{A second, independent grid.} To assess whether the conclusion is specific to
  California, we applied the unchanged methodology to Spain, Portugal and France. Spain and Portugal
  share the MIBEL day-ahead market and are tightly coupled, directly addressing the requirement that
  the regions be electrically coupled to be of interest.
\end{enumerate}
The optimization method itself --- the Mahalanobis--Wasserstein second-order cone program (SOCP) of
Algorithm~2b --- is held fixed throughout; only the feasible set and the region set vary.

\section{Methodology}

\subsection{The scheduling problem}
We schedule compute across $R$ regions over a $T=24$-hour day. The decision variable is
\[
  x \in \R^{R\times T}_{\ge 0}, \qquad x_{r,t} = \text{compute power in region $r$ at hour $t$ [MW]}.
\]
The carbon intensity of each region--hour, $\rho_{r,t}$ (in $\gco$), is uncertain. The carbon cost of a
schedule is linear in $\rho$:
\[
  \langle \rho, x\rangle = \sum_{r=1}^{R}\sum_{t=0}^{T-1} \rho_{r,t}\, x_{r,t}.
\]

\subsection{Distributionally robust objective}
Rather than trust a single forecast, we minimize the worst-case expected cost over a Wasserstein
ambiguity ball $\mathcal{B}_\varepsilon(\hat{\mathbb{P}})$ of radius $\varepsilon\ge 0$ centered on the
empirical distribution $\hat{\mathbb{P}}$ of the daily carbon field:
\begin{equation}
  \min_{x\in\mathcal{X}}\ \sup_{\mathbb{Q}\in\mathcal{B}_\varepsilon(\hat{\mathbb{P}})}\
  \E_{\mathbb{Q}}\!\big[\langle \rho, x\rangle\big].
  \label{eq:dro}
\end{equation}
For a cost linear in $\rho$ and a type-2 Wasserstein ball whose ground metric is the Mahalanobis
distance induced by the empirical covariance $\hat\Sigma$, the inner supremum admits the closed form of
Wasserstein-DRO duality \cite{esfahani2018}: the worst case equals the empirical-mean cost plus
$\varepsilon$ times the Mahalanobis norm of the schedule. Problem~\eqref{eq:dro} reduces to
\begin{equation}
  \boxed{\ \min_{x\in\mathcal{X}}\ \underbrace{\langle \bar\rho, x\rangle}_{\text{mean carbon}}
  \;+\; \underbrace{\varepsilon\,\sqrt{x^{\top}\hat\Sigma\,x}}_{\text{robustness penalty}}\ }
  \label{eq:obj}
\end{equation}
where $\bar\rho\in\R^{R\times T}$ is the empirical mean field and $\hat\Sigma\in\R^{RT\times RT}$ is the
empirical covariance of the flattened daily field $\vecx(\rho)$, using the fixed row-major convention
$\vecx(x)_{rT+t}=x_{r,t}$. The penalty $\sqrt{x^\top\hat\Sigma x}$ is the standard deviation of the
schedule's carbon cost under $\hat\Sigma$; the radius $\varepsilon$ controls the strength of the hedge,
with $\varepsilon=0$ recovering the deterministic problem on $\bar\rho$.

\subsection{Where spatial correlation enters}
Partition $\hat\Sigma$ into $T\times T$ blocks indexed by region pairs,
\begin{equation}
  \hat\Sigma=\begin{pmatrix}\Sigma_{11}&\cdots&\Sigma_{1R}\\ \vdots&\ddots&\vdots\\
  \Sigma_{R1}&\cdots&\Sigma_{RR}\end{pmatrix},
  \qquad \Sigma_{rs}=\Cov\!\big(\rho_{r,\cdot},\,\rho_{s,\cdot}\big)\in\R^{T\times T}.
  \label{eq:blocks}
\end{equation}
The diagonal blocks $\Sigma_{rr}$ are within-region temporal covariances; the off-diagonal blocks
$\Sigma_{rs}$ ($r\ne s$) are exactly the cross-region (spatial) covariance. Through the quadratic form
\begin{equation}
  x^\top\hat\Sigma x \;=\; \sum_{r} x_r^\top \Sigma_{rr}\, x_r \;+\; \sum_{r\ne s} x_r^\top \Sigma_{rs}\, x_s,
  \label{eq:quad}
\end{equation}
the off-diagonal blocks are the \emph{only} channel through which spatial correlation can change the
schedule. The experiment in \S\ref{sec:test} manipulates precisely these blocks.

\subsection{Second-order cone reformulation}
We factor the regularized covariance by its Cholesky decomposition $\hat\Sigma = L L^{\top}$ with $L$
lower triangular. Then $\sqrt{x^{\top}\hat\Sigma x}=\sqrt{x^{\top}LL^{\top}x}=\lVert L^{\top}\vecx(x)\rVert_2$,
so \eqref{eq:obj} becomes
\begin{equation}
  \min_{x\in\mathcal{X}}\ \langle\bar\rho,x\rangle+\varepsilon\,\big\lVert L^{\top}\vecx(x)\big\rVert_{2},
  \label{eq:socp}
\end{equation}
which, via its epigraph $\min\,\langle\bar\rho,x\rangle+\varepsilon s$ subject to
$\lVert L^{\top}\vecx(x)\rVert_2\le s$, is a convex second-order cone program solved to global
optimality (CLARABEL). The orientation $L^{\top}$ is essential: it yields $x^\top(LL^\top)x=x^\top\hat\Sigma x$,
whereas $\lVert L\,\vecx(x)\rVert_2$ would compute the different form $x^\top L^\top L\,x$; the
implementation pins the correct orientation with a post-solve numerical cross-check.

\subsection{The revised feasible set \texorpdfstring{$\mathcal{X}$}{X}}
Every constraint below is linear, so \eqref{eq:socp} remains an SOCP.

\paragraph{C0 --- per-cell ceiling (kept).}
$0\le x_{r,t}\le \bar x_{r,t}$; a site cannot exceed its installed capacity ($\bar x=50$ MW).

\paragraph{C2 --- flexible/inflexible split (kept).}
A fraction $\alpha_r\in[0,1]$ of each region's daily work $W_r$ is inflexible on a fixed intraday shape
$p_{r,t}$ ($\sum_t p_{r,t}=1$); the remainder is flexible:
\begin{equation}
  x_{r,t}=\alpha_r W_r\,p_{r,t}+x^{\mathrm{flex}}_{r,t},\quad x^{\mathrm{flex}}_{r,t}\ge 0,\quad
  \sum_{t} x^{\mathrm{flex}}_{r,t}=(1-\alpha_r)\,W_r.
  \label{eq:c2}
\end{equation}
Here $\alpha$ is the master flexibility dial: lower $\alpha$ leaves more movable work and thus more room
for a spatial effect. We sweep $\alpha\in\{0.30,0.50,0.75\}$ and fix utilization so $W_r=0.80\,\bar x\,T=960$ MWh.

\paragraph{C3 --- inter-hour ramp (kept).}
$\lvert x_{r,t}-x_{r,t-1}\rvert\le\Delta_r$ for $t=1,\dots,T-1$, with $\Delta_r=15$ MW/h. This makes the
naive ``cheapest-hours-first'' schedule infeasible.

\paragraph{C1 --- aggregate power cap (dropped).}
The previous constraint $\sum_r x_{r,t}\le P_{\max}$ is removed; it was a single global knob that
flattened the feasible polytope.

\paragraph{3a --- windowed-demand / deferral deadline (new).}
For a window $[\tau_1,\tau_2]$ and fraction $\gamma$, a $\gamma$-share of each region's flexible work
must clear inside the window:
\begin{equation}
  \sum_{t=\tau_1}^{\tau_2} x^{\mathrm{flex}}_{r,t}\ \ge\ \gamma\,(1-\alpha_r)\,W_r.
  \label{eq:3a}
\end{equation}
This is an aggregate deadline (a floor on flexible work completed by the window's end), not a per-job
service-level agreement. We use the morning window $[0,7]$ with $\gamma=0.20$.

\paragraph{3b --- temperature-coupled thermal/PUE constraint (new).}
Usable IT power is limited by effective (wall-plug) power, the IT power inflated by the facility PUE,
which rises with outdoor temperature through a piecewise-linear ``hockey stick'':
\begin{equation}
  \mathrm{PUE}(T)=\mathrm{PUE}_0+\kappa\,\max\!\big(T-T_{\mathrm{set}},\,0\big),\qquad
  \mathrm{PUE}\!\big(T^{\mathrm{air}}_{r,t}\big)\,x_{r,t}\ \le\ \bar P.
  \label{eq:3b}
\end{equation}
Because the air temperature $T^{\mathrm{air}}_{r,t}$ is data (a fixed field, never re-estimated per fold
nor read from the test set), $\mathrm{PUE}(\cdot)$ is a known constant per cell and \eqref{eq:3b} is
linear --- it tightens the ceiling to $\bar P/\mathrm{PUE}(T^{\mathrm{air}}_{r,t})$ in hot hours. We use
$\mathrm{PUE}_0=1.10$, $\kappa=0.015\,\mathrm{/}\degC$, $\bar P=57$ MW, and economizer set-point
$T_{\mathrm{set}}=20\degC$ (California) or $14\degC$ (the milder Iberian field; see \S\ref{sec:reliable}).

\subsection{The falsification test: joint versus shuffled covariance}\label{sec:test}
To isolate the value of \emph{spatial} correlation specifically, we feed the scheduler two covariances
that differ in exactly one respect. The joint covariance $\hat\Sigma^{\mathrm{joint}}$ is the full
empirical covariance \eqref{eq:blocks}. The shuffled covariance $\hat\Sigma^{\mathrm{shuf}}$ zeros all
cross-region blocks while preserving the within-region blocks:
\begin{equation}
  \big(\hat\Sigma^{\mathrm{shuf}}\big)_{rs}=
  \begin{cases}\Sigma_{rr}, & r=s,\\ \mathbf{0}_{T\times T}, & r\ne s.\end{cases}
  \label{eq:shuf}
\end{equation}
This preserves each region's marginal distribution and within-region temporal autocorrelation, and
destroys only cross-region dependence. Hence comparing the scheduler under $\hat\Sigma^{\mathrm{joint}}$
versus $\hat\Sigma^{\mathrm{shuf}}$ isolates exactly the value of spatial correlation. If that structure
is useful, the joint-fitted schedule should achieve lower out-of-sample tail cost. The test is a
\emph{falsification}: it is built to detect a spatial benefit if one exists, so a null is informative.

\subsection{Evaluation protocol}
\paragraph{Risk metric.} A schedule $x$ is scored on each held-out test day $i$ by its realized
emissions $e_i=\langle\rho_i,x\rangle$. The objective is the Conditional Value-at-Risk at level
$\beta=0.95$ --- the mean of the worst $5\%$ of days:
\begin{equation}
  \CVaR_{\beta}\big(\{e_i\}_{i=1}^{N}\big)=\frac{1}{\lceil(1-\beta)N\rceil}\sum_{i\in\mathcal{T}_\beta}e_i,
  \qquad \mathcal{T}_\beta=\{\text{indices of the largest }\lceil(1-\beta)N\rceil\ e_i\}.
  \label{eq:cvar}
\end{equation}
We report the \emph{spatial gap}
$\ \mathrm{gap}=\CVaR_{0.95}(\mathrm{shuf})-\CVaR_{0.95}(\mathrm{joint})$,
so a positive gap means the joint covariance helped; we express it as a percentage of the shuffled CVaR
so the two grids are comparable.

\paragraph{Selecting $\varepsilon$.} For each arm separately, $\varepsilon$ is chosen by 5-fold
\emph{blocked} (contiguous, non-shuffled) time-series cross-validation on 2021--2024 over the grid
$\{0,0.1,1,10,100,1000\}$, minimizing mean validation CVaR; the selected $\varepsilon^\star$ is refit on
the full training set and evaluated once on the untouched 2025 test set.

\paragraph{Uncertainty.} A $95\%$ confidence interval on the gap is obtained by a paired bootstrap:
test-day indices are resampled with replacement (the same indices for both arms), the gap recomputed on
each of $1000$ resamples, and the $2.5/97.5$ percentiles reported; a gap is \emph{detectable} only if
the interval excludes zero.

\paragraph{Pre-registration.} For every experiment we (i) ran a dry-run gate performing the full
cross-validation without touching 2025; (ii) commit-locked the configuration in version control before
any test-set evaluation; (iii) read the test set exactly once. This removes the analyst degrees of
freedom that could bias the outcome.

\section{Experimental design}
The joint-versus-shuffled comparison is run across three nested constraint regimes so that the value of
spatial correlation is reported as a function of operational tightness (Table~\ref{tab:regimes}). Two
region sets are studied with \emph{identical} method, metric, cross-validation, bootstrap and table
format; only what must differ does.

\begin{table}[h]\centering
\caption{The three nested constraint regimes (each adds to the previous).}
\label{tab:regimes}
\begin{tabular}{@{}lll@{}}
\toprule
Regime & Active constraints & Role \\
\midrule
R3 (reference) & C0 + C2 + C3 & post-cap baseline (anchor) \\
R1 (lean)      & R3 + deadline (3a) & headline regime; no weather assumptions \\
R2 (full)      & R1 + thermal (3b) & tightest, most assumption-laden \\
\bottomrule
\end{tabular}
\end{table}

\paragraph{Data.} Carbon intensity is hourly life-cycle intensity from Electricity Maps (academic tier),
2021--2025. Temperature is hourly 2\,m air temperature from the Open-Meteo ERA5 archive, one
load-weighted point per zone (California: Fresno, Sacramento, Los Angeles, Las Vegas; Iberia: Madrid,
Lisbon, Paris). Both panels are built on a common local clock, train on 2021--2024 and test on 2025.

\section{Results}

\subsection{Case A --- California / Nevada ($R=4$)}
Zones CISO, BANC, LDWP and NEVP. Test-set $\CVaR_{0.95}\approx 1.59\times 10^{6}\ \mathrm{gCO_2}$.
Table~\ref{tab:ca} reports the spatial gap. The largest benefit anywhere is $+0.058\%$; the cells flagged
detectable at $\alpha=0.75$ are economically meaningless ($0.0017\%$, near-identical schedules).

\begin{table}[h]\centering
\caption{California: spatial gap $\CVaR_{0.95}(\mathrm{shuf})-\CVaR_{0.95}(\mathrm{joint})$ by regime and $\alpha$. Positive favors the joint covariance.}
\label{tab:ca}
\begin{tabular}{@{}llrrl@{}}
\toprule
Regime & $\alpha$ & gap (\%) & $95\%$ CI (gCO$_2$) & verdict \\
\midrule
R3 & 0.30 & $+0.0583$ & $[\,+211,\ +1942\,]$ & detectable \\
R3 & 0.50 & $+0.0357$ & $[\,-532,\ +2177\,]$ & null \\
R3 & 0.75 & $-0.0122$ & $[\,-581,\ +488\,]$ & null \\
R1 & 0.30 & $+0.0164$ & $[\,-118,\ +831\,]$ & null \\
R1 & 0.50 & $+0.0271$ & $[\,-359,\ +1577\,]$ & null \\
R1 & 0.75 & $+0.0017$ & $[\,+13,\ +41\,]$ & detectable$^\ast$ \\
R2 & 0.30 & $+0.0185$ & $[\,+34,\ +693\,]$ & detectable \\
R2 & 0.50 & $+0.0262$ & $[\,-267,\ +1366\,]$ & null \\
R2 & 0.75 & $+0.0017$ & $[\,+13,\ +41\,]$ & detectable$^\ast$ \\
\bottomrule
\end{tabular}
\\[2pt]{\footnotesize $^\ast$ Statistically detectable but $0.0017\%$ of CVaR --- practically null.}
\end{table}

\subsection{Case B --- Iberia + France ($R=3$)}
Zones ES, PT, FR. Test-set $\CVaR_{0.95}\approx 0.40\times 10^{6}\ \mathrm{gCO_2}$ (lower than
California because French intensity is low: mean $\approx 50$ versus ES/PT $\approx 160\ \gco$).
Table~\ref{tab:ib} reports the gap; eight of nine cells are practically null.

\begin{table}[h]\centering
\caption{Iberia + France: spatial gap $\CVaR_{0.95}$ by regime and $\alpha$.}
\label{tab:ib}
\begin{tabular}{@{}llrrl@{}}
\toprule
Regime & $\alpha$ & gap (\%) & $95\%$ CI (gCO$_2$) & verdict \\
\midrule
R3 & 0.30 & $+0.0000$ & $\varepsilon^\star{=}0$, schedules identical & null \\
R3 & 0.50 & $-0.0000$ & $\varepsilon^\star{=}0.1$, $\approx$ identical & null \\
R3 & 0.75 & $+0.0210$ & $[\,-108,\ +226\,]$ & null \\
R1 & 0.30 & $-0.0017$ & $[\,-843,\ +678\,]$ & null \\
R1 & 0.50 & $+0.0151$ & $[\,-18,\ +129\,]$ & null \\
R1 & 0.75 & $+0.3211$ & $[\,+746,\ +1870\,]$ & detectable$^\dagger$ \\
R2 & 0.30 & $-0.0075$ & $[\,-563,\ +495\,]$ & null \\
R2 & 0.50 & $+0.0000$ & $[\,-0.004,\ +0.013\,]$ & null \\
R2 & 0.75 & $-0.0155$ & $[\,-152,\ +19\,]$ & null \\
\bottomrule
\end{tabular}
\\[2pt]{\footnotesize $^\dagger$ A cross-validation selection artifact; see \S\ref{sec:artifact}.}
\end{table}

\subsection{The apparent exception and why it does not survive}\label{sec:artifact}
The Iberian R1/$\alpha{=}0.75$ cell shows a raw $+0.32\%$ gap with a confidence interval excluding zero
--- the single largest raw gap in either study. It arises \emph{only} because cross-validation hands the
two arms different radii ($\varepsilon^\star_{\mathrm{joint}}=1$ but
$\varepsilon^\star_{\mathrm{shuf}}=10$), i.e.\ it compares two schedules at different penalties rather
than measuring an informational advantage. When the comparison is made clean --- residual-space
estimation, under which both arms select $\varepsilon^\star=1$ --- the sign flips
(Table~\ref{tab:artifact}).

\begin{table}[h]\centering
\caption{The R1/$\alpha{=}0.75$ cell across estimators. The raw effect reverses under clean (equal-$\varepsilon^\star$) comparison.}
\label{tab:artifact}
\begin{tabular}{@{}lccl@{}}
\toprule
Estimator & $\varepsilon^\star$ joint/shuf & gap (\%) & clean comparison? \\
\midrule
Sample covariance      & $1/10$ & $+0.321$ & no --- $\varepsilon^\star$ diverge \\
Ledoit--Wolf shrinkage & $1/10$ & $+0.321$ & no --- $\varepsilon^\star$ diverge \\
Seasonal residual      & $1/1$  & $-0.032$ & yes \\
AR(1) residual         & $1/1$  & $-0.037$ & yes \\
\bottomrule
\end{tabular}
\end{table}

Under both clean estimators the joint covariance is marginally \emph{worse}, and they agree in sign; our
pre-registered ``seasonal and AR(1) must agree'' rule therefore records no beneficial spatial effect. The
apparent exception is a selection artifact, caught by the robustness machinery designed to catch it.

\paragraph{Verdict.} In both grids the gap is a few hundredths of a percent of CVaR essentially
everywhere; neither shows an economically meaningful, robust spatial benefit. \textbf{Iberia replicates
the California null; the more tightly coupled MIBEL grid did not surface spatial value.}

\section{Reliability of the result}\label{sec:reliable}

\paragraph{Substantive, not mechanical (the Goldilocks check).} A null is meaningful only if the
optimizer \emph{could} have responded to spatial structure. We verify the feasible set is not frozen by
measuring how far the deterministic optimum reallocates as $\varepsilon$ grows,
$\mathrm{move}(\varepsilon)=\lVert x(\varepsilon)-x(0)\rVert_1$ as a fraction of total work. Across all
nine regime$\times\alpha$ cells the optimizer retains substantial freedom --- $24$--$32\%$ in California
and $9.5$--$18\%$ in Iberia --- so the null reflects a genuine absence of exploitable spatial structure.
By design every constraint is active but not dominant (in R2 the thermal cap binds $\approx 22$--$26\%$
of cells; the deadline binds for most regions), present and realistic without removing the optimizer's
freedom.

\paragraph{Robustness battery.} (i) \emph{Ledoit--Wolf shrinkage} \cite{ledoit2004}: re-running with the
shrinkage covariance reproduces every gap to $\approx 4$ significant figures, so the result is not an
artifact of a noisy sample covariance. (ii) \emph{Residual-space}: re-estimating the covariance from the
forecast residuals of a seasonal and an AR(1) baseline (keeping the mean field and scored emissions on
raw intensity), with a pre-registered agreement rule, yields no beneficial spatial effect in either
grid.

\paragraph{Replication across independent grids.} Two grids differing in geography, climate, generation
mix and coupling structure (CAISO/Nevada and MIBEL/France) give the same conclusion. A replicated null
under a falsification test is substantially stronger evidence than any single region.

\paragraph{Reproducibility.} Every value here is read from committed result files; the pipeline carries
$149$ unit tests, and the configuration was version-locked before the test set was scored. The single
re-calibrated parameter across the two grids was the thermal set-point ($20\to 14\degC$), because the
milder Iberian annual-mean field would otherwise leave constraint~\eqref{eq:3b} inactive; the headline
regime R1 carries no thermal assumption, so the conclusion does not depend on this choice.

\section{Interpretation: why this is the expected outcome}
High \emph{marginal} correlation does not imply exploitable \emph{robust-scheduling} value. In Iberia the
Spain--Portugal carbon-intensity correlation is $\approx 0.85$, higher than any California pair
($0.43$--$0.77$), yet the joint covariance buys nothing over the block-diagonal one. Two mechanisms
explain this. First, cross-validation selects small radii ($\varepsilon^\star\in\{0,\dots,1\}$), so the
robustness term in \eqref{eq:socp} is a minor correction on top of the mean-cost term shared identically
by both arms; the schedule is driven mostly by where carbon is cheap on average. Second, in the
quadratic form \eqref{eq:quad} the within-region (diagonal) terms dominate for a single fixed schedule
concentrated in each region's cheap hours, so zeroing the off-diagonal blocks $\Sigma_{rs}$ changes the
penalty --- and hence the optimal schedule and its tail cost --- by very little. Spatial structure would
have the most leverage if the operator could continuously reallocate \emph{between} regions in response
to a realized joint shock; the present design commits to one representative schedule, exactly the
setting in which spatial correlation has the least purchase.

\section{Limitations and next steps}
\begin{itemize}[leftmargin=1.4em,itemsep=2pt]
  \item \textbf{One representative schedule.} A rolling day-ahead formulation that re-optimizes each day
  against updated forecasts is the most likely route to a genuine spatial benefit and the highest-value
  next experiment.
  \item \textbf{Common-clock convention.} Portugal (WET) is scheduled on the Spain/France (CET) clock, a
  one-hour offset; this cannot bias the joint-versus-shuffled comparison (it perturbs both arms
  identically) and is documented.
  \item \textbf{Reporting clarity.} The detectability flag will be supplemented with a practical-
  significance threshold so trivially small but CI-tight cells are not misread.
  \item \textbf{Publication track.} Given the replicated null, the planned vine-copula extension may
  pivot toward characterizing \emph{when} spatial structure matters (metric, ambiguity set, or the
  adaptive setting above) rather than presuming it does.
\end{itemize}

\section{Conclusion}
On a defensible feasible set, under a falsification test evaluated by out-of-sample CVaR with strict
pre-registration, spatial correlation of carbon intensity provides no robust scheduling value --- a
result that replicates across two independent grids and survives a shrinkage and residual-space
robustness battery, and which we show to be substantive rather than a mechanical artifact. The
contribution this period is a clean, replicated negative result with a mechanistic explanation, and a
clear next experiment (adaptive re-optimization) under which a positive effect, if any exists, would be
most likely to appear.

\begin{thebibliography}{9}
\bibitem{esfahani2018} P.~Mohajerin Esfahani and D.~Kuhn. Data-driven distributionally robust
optimization using the Wasserstein metric. \emph{Mathematical Programming}, 171:115--166, 2018.
\bibitem{hall2024} G.~Hall et al. Distributionally robust carbon-aware scheduling with virtual capacity
curves. arXiv:2410.21510, 2024.
\bibitem{ledoit2004} O.~Ledoit and M.~Wolf. A well-conditioned estimator for large-dimensional
covariance matrices. \emph{Journal of Multivariate Analysis}, 88(2):365--411, 2004.
\bibitem{radovanovic2022} A.~Radovanovi\'c et al. Carbon-aware computing for data centers.
\emph{IEEE Transactions on Power Systems}, 38(2):1270--1280, 2022.
\end{thebibliography}

\appendix
\section{Reproducibility}
Repository \texttt{ma731/Capstone-in-Operations-Research}, branch \texttt{main}. Result files:
\texttt{results/taskA\_regimes\_2026-06-02.csv} (California) and
\texttt{results/es\_pt\_fr\_regimes\_2026-06-05*.csv} (Iberia + France, including the
\texttt{\_lw}, \texttt{\_seasonal} and \texttt{\_ar1} robustness variants). Detailed working notes:
\texttt{docs/progress\_note\_v13.md} and \texttt{docs/progress\_note\_v14.md}. Mathematics:
\texttt{thesis/full\_formulation.md} and \texttt{thesis/constraints\_explained.md}. The optimization
core is \texttt{src/models/algorithm\_2b\_mahalanobis.py}; the locked experiment drivers are
\texttt{scripts/run\_shuffled\_marginals\_taskA\_experiment.py} and
\texttt{scripts/run\_shuffled\_marginals\_es\_pt\_fr\_experiment.py}.

\end{document}
